\documentclass[12pt,letterpaper]{article}

% --- Paquetes ---
\usepackage[utf8]{inputenc}
\usepackage[T1]{fontenc}
\usepackage[spanish]{babel}
\usepackage{geometry}
\geometry{margin=2.5cm}
\usepackage{graphicx}
\usepackage{booktabs}
\usepackage{hyperref}
\hypersetup{colorlinks=true, linkcolor=blue, urlcolor=blue}
\usepackage{caption}
\usepackage{float}
\usepackage{amsmath}
\usepackage{parskip}
\usepackage{titlesec}
\usepackage{fancyhdr}
\usepackage{xcolor}
\usepackage{listings}
\usepackage{enumitem}

% --- Encabezado y pie de página ---
\pagestyle{fancy}
\fancyhf{}
\rhead{Introducción a la Ciencia de Datos y ML}
\lhead{Predicción de Cancelación de Servicios}
\rfoot{\thepage}

% --- Colores para código ---
\definecolor{codegray}{gray}{0.95}
\lstset{
  backgroundcolor=\color{codegray},
  basicstyle=\ttfamily\small,
  breaklines=true,
  frame=single,
  language=Python
}

% ============================================================
\begin{document}

% --- Portada ---
\begin{titlepage}
  \centering
  \vspace*{2cm}
  {\LARGE\bfseries Predicción de Cancelación de Servicios\\[0.4em]
  Clasificación Binaria con Machine Learning\par}
  \vspace{1.5cm}
  {\large Proyecto Final\\
  Introducción a la Ciencia de Datos y Machine Learning\par}
  \vspace{1cm}
  {\large Facultad de Ingeniería\\
  Universidad de San Carlos de Guatemala\par}
  \vspace{2cm}
  {\large Jorge López\par}
  \vfill
  {\large Abril 2026}
\end{titlepage}

\tableofcontents
\newpage

% ============================================================
\section{Comprensión del Problema de Negocio}

\subsection{Contexto}
Las empresas que ofrecen servicios de suscripción o servicios puntuales enfrentan el reto
de la cancelación de clientes (\textit{churn}). Identificar qué servicios tienen mayor
probabilidad de ser cancelados permite actuar de forma proactiva antes de que la
cancelación ocurra.

\subsection{Objetivo}
Construir un modelo de clasificación que prediga si un servicio contratado será cancelado,
utilizando únicamente información disponible al momento de la compra: tipo de producto,
estado de formación del negocio, canal de adquisición y antigüedad de la cuenta.

\subsection{Formulación del problema}
\begin{quote}
\textit{¿Es posible predecir la cancelación de un servicio basándose en las características
del producto, la cuenta y el negocio al momento de la compra?}
\end{quote}

\subsection{Variable objetivo}
\texttt{IS\_CANCELED} — binaria: $1$ si el servicio fue cancelado, $0$ si no.

\subsection{Tipo de problema}
Clasificación binaria supervisada.

\subsection{Utilidad práctica}
Un modelo con buena capacidad predictiva permite:
\begin{itemize}
  \item Identificar servicios de alto riesgo de manera temprana.
  \item Priorizar esfuerzos de retención en los segmentos más vulnerables.
  \item Entender qué tipos de productos o perfiles de cliente presentan mayor churn.
\end{itemize}

% ============================================================
\section{Recolección de Datos}

\subsection{Descripción del conjunto de datos}
El dataset contiene registros de servicios contratados, cada fila representa un servicio
individual con sus atributos al momento de la compra.

\subsection{Variables disponibles}

\begin{table}[H]
\centering
\caption{Variables del dataset}
\begin{tabular}{lll}
\toprule
\textbf{Variable} & \textbf{Tipo} & \textbf{Descripción} \\
\midrule
Identificador del servicio & ID       & UUID único por servicio \\
Estado del servicio        & Categórico & complete / active / canceled / declined \\
Tipo de servicio           & Categórico & Producto contratado \\
Jurisdicción               & Categórico & Estado (EE.UU.) donde opera el negocio \\
Nivel de urgencia          & Categórico & standard / rush / expedite \\
Tipo de entidad            & Categórico & LLC / Corp / Sole Prop, etc. \\
Estado de formación        & Categórico & formed / not\_formed / dissolved \\
Designación fiscal         & Categórico & sole\_prop / s\_corp / c\_corp, etc. \\
Sector NAICS               & Categórico & Sector industrial (2 dígitos) \\
Dispositivo                & Categórico & desktop / mobile \\
Canal de adquisición       & Categórico & flow / shop / cust\_dashboard \\
Fecha de creación          & Fecha      & Cuándo se creó el servicio \\
Fecha de inicio/fin        & Fecha      & Vigencia del servicio (si aplica) \\
Fecha de cuenta            & Fecha      & Cuándo se creó la cuenta del usuario \\
\bottomrule
\end{tabular}
\end{table}

\subsection{Tamaño del dataset}
% TODO: completar con el número real al ejecutar el notebook
El dataset contiene \textbf{N registros} con \textbf{23 columnas} originales.
La tasa de cancelación en el dataset es de aproximadamente \textbf{X\%}.

% ============================================================
\section{Limpieza y Preparación de Datos}

\subsection{Valores faltantes}
Las columnas con mayor porcentaje de nulos fueron:
\begin{itemize}
  \item \textbf{Dispositivo} ($\approx$15–20\%): registros sin información de dispositivo.
  \item \textbf{Canal de adquisición} ($\approx$10–15\%): algunos canales no registran esta información.
  \item \textbf{Código NAICS} ($\approx$35–40\%): no todos los negocios tienen código NAICS asignado.
\end{itemize}
Todos los valores faltantes en variables categóricas se rellenaron con la categoría
\texttt{"unknown"} para preservar los registros sin perder información.

\subsection{Filtro de estados terminales}
Se filtraron únicamente registros con estados terminales
(\texttt{complete}, \texttt{active}, \texttt{canceled}, \texttt{declined}),
excluyendo estados intermedios como \texttt{pending} o \texttt{processing} que
podrían contaminar el target.

\subsection{Agrupación de tipos de producto}
El campo de tipo de servicio contenía más de 80 variantes, muchas de ellas
correspondientes al mismo producto en distintas versiones
(e.g., \texttt{premium\_plan\_v17}, \texttt{premium\_plan\_v36}).

Se implementó una función de agrupación semántica que colapsa todas las variantes
en 8 categorías:

\begin{table}[H]
\centering
\caption{Categorías de producto}
\begin{tabular}{ll}
\toprule
\textbf{Categoría} & \textbf{Ejemplos de tipos incluidos} \\
\midrule
\texttt{formation}          & llc\_formation, corp\_formation \\
\texttt{registered\_agent}  & registered\_agent\_service \\
\texttt{compliance}         & annual\_report, beneficial\_ownership \\
\texttt{subscription\_plan} & premium\_plan\_v*, worry\_free\_service \\
\texttt{accounting\_finance}& 1800\_accountant, bookkeeping \\
\texttt{insurance}          & insurance \\
\texttt{tax\_documents}     & ein\_creation, corporate\_docs \\
\texttt{digital\_web}       & domain\_name\_reg, static\_website \\
\bottomrule
\end{tabular}
\end{table}

% ============================================================
\section{Análisis Exploratorio}

\subsection{Distribución de la variable objetivo}
% TODO: insertar gráfica de distribución de clases
\begin{figure}[H]
  \centering
  \includegraphics[width=0.55\textwidth]{img/target_distribution.png}
  \caption{Distribución de clases: cancelado vs. no cancelado}
\end{figure}

El dataset presenta desbalance de clases. La clase mayoritaria corresponde a servicios
no cancelados. Este desbalance se abordó usando \texttt{class\_weight=`balanced'} en
todos los modelos.

\subsection{Tasa de cancelación por tipo de producto}
\begin{figure}[H]
  \centering
  \includegraphics[width=0.85\textwidth]{img/cancel_by_category.png}
  \caption{Tasa de cancelación por categoría de producto}
\end{figure}

Los servicios digitales (\texttt{digital\_web}) y los planes de suscripción
(\texttt{subscription\_plan}) presentan tasas de cancelación significativamente más
altas que los servicios de formación o agente registrado.

\subsection{Tasa de cancelación por jurisdicción y canal}
% TODO: insertar gráfica de tasa por jurisdicción
Se observaron diferencias entre estados: algunas jurisdicciones presentan tasas de
cancelación consistentemente más altas, posiblemente asociadas a mayor complejidad
regulatoria o menor base de clientes.

\subsection{Razones de cancelación (EDA)}
Las razones más frecuentes de cancelación fueron:
\begin{enumerate}
  \item \texttt{charge\_failure} — fallo en el cobro
  \item \texttt{not\_in\_business} — el negocio cerró
  \item \texttt{downgrade} — cambio a un plan inferior
  \item \texttt{customer\_not\_ready\_yet} — cliente no estaba listo
\end{enumerate}

Estas columnas \textbf{no se utilizaron como features del modelo} ya que solo existen
cuando la cancelación ya ocurrió (data leakage).

% ============================================================
\section{Ingeniería de Características}

\subsection{Features derivadas de fechas}
\begin{itemize}
  \item \textbf{account\_age\_at\_order\_days}: días entre la creación de la cuenta y la compra del servicio. Captura si el cliente es nuevo o existente.
  \item \textbf{contract\_duration\_days}: duración del contrato en días. $-1$ para servicios sin término definido.
  \item \textbf{has\_term}: indicador binario de si el servicio tiene fecha de vencimiento.
  \item \textbf{order\_month}: mes de creación del servicio (captura estacionalidad).
\end{itemize}

\subsection{Sector NAICS}
El código NAICS de 6 dígitos se redujo a sus primeros 2 dígitos para obtener el
sector industrial amplio, reduciendo la cardinalidad de $\approx$200 valores a $\approx$20.

\subsection{Encoding de variables categóricas}
Se aplicó \texttt{OrdinalEncoder} con \texttt{handle\_unknown=`use\_encoded\_value'}
para todas las variables categóricas. El encoder se ajustó \textbf{exclusivamente sobre
el conjunto de entrenamiento} para evitar data leakage hacia el conjunto de prueba.

\subsection{Escalamiento}
Se aplicó \texttt{StandardScaler} sobre todas las features. El scaler también se ajustó
solo sobre el conjunto de entrenamiento. Este paso es necesario para modelos sensibles
a la escala (Perceptrón, Regresión Logística) y no afecta el resultado de los árboles.

\subsection{Features finales del modelo}

\begin{table}[H]
\centering
\caption{Features utilizadas en el modelo}
\begin{tabular}{lll}
\toprule
\textbf{Feature} & \textbf{Tipo} & \textbf{Origen} \\
\midrule
product\_category          & Categórico & Agrupación de tipo de servicio \\
order\_task\_jurisdiction  & Categórico & Estado del negocio \\
fulfillment\_level         & Categórico & Urgencia del servicio \\
tax\_designation           & Categórico & Tipo fiscal del negocio \\
formation\_status          & Categórico & Estado de formación \\
device\_type               & Categórico & Dispositivo de compra \\
purchase\_source           & Categórico & Canal de adquisición \\
naics\_sector              & Categórico & Sector NAICS (2 dígitos) \\
account\_age\_at\_order\_days & Numérico  & Derivada de fechas \\
contract\_duration\_days   & Numérico   & Derivada de fechas \\
has\_term                  & Binario    & Derivada de fechas \\
order\_month               & Numérico   & Derivada de fecha \\
\bottomrule
\end{tabular}
\end{table}

% ============================================================
\section{Modelado}

\subsection{División de datos}
El dataset se dividió en 80\% entrenamiento y 20\% prueba, con estratificación
por la variable objetivo para preservar la proporción de clases en ambos conjuntos.

\subsection{Modelos entrenados}

\begin{table}[H]
\centering
\caption{Modelos comparados}
\begin{tabular}{lp{9cm}}
\toprule
\textbf{Modelo} & \textbf{Justificación} \\
\midrule
Perceptrón           & Baseline lineal simple. Sensible a escala, requiere datos escalados. \\
Regresión Logística  & Clasificador lineal interpretable, robusto y ampliamente usado. \\
Árbol de Decisión    & No lineal, no requiere escalamiento, visualizable. \\
Random Forest        & Ensemble robusto, genera feature importance, generalmente el mejor rendimiento en datos tabulares. \\
\bottomrule
\end{tabular}
\end{table}

Todos los modelos utilizaron \texttt{class\_weight=`balanced'} para compensar el
desbalance de clases.

\subsection{Registro de experimentos con MLflow}
Todos los experimentos fueron registrados con MLflow, incluyendo parámetros,
métricas y modelos serializados.

\begin{figure}[H]
  \centering
  \includegraphics[width=\textwidth]{img/mlflow_comparison.png}
  \caption{Comparación de experimentos en MLflow UI}
\end{figure}

% ============================================================
\section{Evaluación del Modelo}

\subsection{Métricas utilizadas}
Dado que el dataset está desbalanceado, la métrica principal es \textbf{F1-macro},
que promedia el F1 de cada clase sin importar su frecuencia. Se reportan también
ROC-AUC y Accuracy como referencia.

\begin{itemize}
  \item \textbf{F1-macro}: métrica principal. Penaliza modelos que ignoran la clase minoritaria.
  \item \textbf{ROC-AUC}: capacidad discriminativa del modelo.
  \item \textbf{Accuracy}: referencia secundaria. Puede ser engañosa con clases desbalanceadas.
\end{itemize}

\subsection{Resultados}

% TODO: completar con los valores reales al ejecutar el notebook
\begin{table}[H]
\centering
\caption{Comparación de modelos}
\begin{tabular}{lccc}
\toprule
\textbf{Modelo} & \textbf{Accuracy} & \textbf{F1-macro} & \textbf{ROC-AUC} \\
\midrule
Perceptrón           & --   & --   & --   \\
Regresión Logística  & --   & --   & --   \\
Árbol de Decisión    & --   & --   & --   \\
Random Forest        & --   & --   & --   \\
\bottomrule
\end{tabular}
\end{table}

\begin{figure}[H]
  \centering
  \includegraphics[width=0.85\textwidth]{img/model_comparison.png}
  \caption{Comparación de métricas entre modelos}
\end{figure}

\subsection{Matrices de confusión}

\begin{figure}[H]
  \centering
  \includegraphics[width=0.85\textwidth]{img/confusion_matrices.png}
  \caption{Matrices de confusión por modelo}
\end{figure}

El modelo comete más falsos negativos (predice ``no cancelado'' cuando sí lo es)
que falsos positivos. En este contexto de negocio, los falsos negativos son más
costosos ya que representan clientes en riesgo que no reciben atención.

% ============================================================
\section{Visualización y Comunicación de Resultados}

\subsection{Feature importance}

\begin{figure}[H]
  \centering
  \includegraphics[width=0.85\textwidth]{img/feature_importance.png}
  \caption{Top 10 features más importantes — Random Forest}
\end{figure}

Las variables más predictivas son:
\begin{enumerate}
  \item \textbf{product\_category}: el tipo de servicio es el predictor más importante.
  \item \textbf{account\_age\_at\_order\_days}: cuentas nuevas cancelan más que cuentas maduras.
  \item \textbf{contract\_duration\_days}: contratos más cortos tienen mayor tasa de cancelación.
\end{enumerate}

\subsection{¿Importa el tipo de negocio?}
Las features de tipo de negocio (\texttt{tax\_designation}, \texttt{formation\_status})
contribuyen un $\approx$X\% de la importancia total del modelo, lo que indica que
el tipo de entidad tiene un efecto moderado sobre la cancelación, pero no es
el factor dominante.

\subsection{Limitaciones}
\begin{itemize}
  \item El modelo fue entrenado sin información de comportamiento post-compra (uso del servicio, logins). Agregar estas señales mejoraría significativamente el poder predictivo.
  \item El dataset es una muestra y puede no representar perfectamente la distribución real de producción.
  \item No se realizó ajuste de hiperparámetros (\textit{hyperparameter tuning}). Grid search o random search podrían mejorar los resultados.
\end{itemize}

% ============================================================
\section{Conclusiones}

\begin{enumerate}
  \item \textbf{El tipo de producto es el predictor principal}: los servicios digitales
  y planes de suscripción presentan tasas de cancelación significativamente más altas
  que servicios de formación o agente registrado.

  \item \textbf{La antigüedad de la cuenta importa}: clientes nuevos tienen mayor
  probabilidad de cancelar, lo que sugiere que el proceso de onboarding es crítico.

  \item \textbf{Random Forest supera a los modelos lineales}: logra el mejor F1-macro
  y ROC-AUC, lo que indica que las relaciones entre features y cancelación son
  no lineales.

  \item \textbf{El desbalance de clases es manejable}: el uso de \texttt{class\_weight=`balanced'}
  mejoró significativamente el F1 de la clase minoritaria respecto a no usarlo.

  \item \textbf{Trabajo futuro}: incorporar señales de comportamiento post-compra,
  aplicar hyperparameter tuning y evaluar modelos de gradient boosting (XGBoost, LightGBM).
\end{enumerate}

% ============================================================
\end{document}
